\chapter{\enquote{Why Did You Use This?}}
\label{ch:why}

Interviewers ask \enquote{why X?} and, almost always, \enquote{why not Y?}. Each answer below gives the reason recorded in the repository, with its source. Where the repository records no reason, the answer is marked \tagI and says so. Answers are written to be spoken.

\section{Approach and data}

\paragraph{Why retrieval-augmented generation rather than fine-tuning?} The goal was answers that can be proven to come from the books, with citations and the ability to abstain; the README states the idea as \enquote{a RAG system should be able to prove its answer is in the source, and abstain when it isn't.} \emph{Why not fine-tuning?} The repository never considers it. \tagI Fine-tuning stores knowledge in weights, so it cannot point to a page, and changing the corpus would mean retraining.

\paragraph{Why these two textbooks?} They are the only CC BY biology titles OpenStax publishes, checked per edition, because second editions moved to a non-commercial licence \ev{src/vidyarag/ingest/corpus.py:16-26}. They also overlap only slightly, which allows genuine cross-book multi-hop questions without the redundancy that would make context precision ambiguous. \emph{Why not the second editions?} Their NonCommercial-ShareAlike licence would restrict reuse of an MIT-licensed repository.

\paragraph{Why parse the PDF outline instead of guessing headings from fonts?} The outline gives real chapter, section and page entries; font-size heuristics \enquote{misfire on pull quotes, figure captions, and section headers that happen to wrap, and\ldots fail silently} \ev{src/vidyarag/ingest/parse.py:8-11}. \emph{Why PDFs rather than the HTML editions?} Not recorded. \tagI The PDFs carry the printed page numbers students can check against a paper copy, which the project treats as essential.

\paragraph{Why 512-token chunks with 64 tokens of overlap?} 512 is the embedding model's hard limit, and it is about a paragraph and a half, which holds a self-contained concept while five chunks stay around 2,600 tokens. The configuration notes the trade-off: 256 fragments definitions, 1,024 dilutes the embedding \ev{config/default.yaml:43-50}. It was not swept.

\paragraph{Why count tokens with the embedding model's tokenizer?} Because a different tokenizer undercounts by about 7\%, so \enquote{512-token} chunks lost their tails at embedding time with no error --- a bug the project hit and fixed \ev{git 9ccd6c2}.

\section{Retrieval}

\paragraph{Why BGE embeddings through fastembed?} They run locally on CPU through ONNX Runtime, need no key and cost nothing, so \enquote{retrieval has no runtime dependency on any third party}, and anyone can build the index without credentials \ev{docs/DESIGN.md}. \emph{Why not OpenAI embeddings?} They needed funded billing. \emph{Why not sentence-transformers?} PyTorch would make the image about 2.5\,GB instead of about 400\,MB.

\paragraph{Why Qdrant?} One client API covers an in-process index, a local server and Qdrant Cloud, so tests, the demo and development differ only in configuration \ev{src/vidyarag/store/client.py:3-11}. \emph{Why not FAISS or Chroma?} Not recorded. \tagI Qdrant gives named vectors, payload filtering and the same code in all three deployment modes.

\paragraph{Why ship an embedded index inside the demo?} Free Qdrant Cloud clusters are suspended after about a week idle and deleted after about four, which would break a portfolio link months later \ev{docs/DESIGN.md}. The cost is exact search, comfortable at 3,608 points and not beyond about 20,000.

\paragraph{Why deterministic point ids?} A \code{uuid5} of the chunk id means re-running ingestion overwrites points instead of duplicating them, so an interrupted run can simply be restarted \ev{src/vidyarag/store/collection.py:3-7}.

\paragraph{Why retrieve 20 and keep 5?} A reranker can only help with a pool larger than the final selection; the configuration calls 4$\times$ a latency and quality sweet spot, though no sweep was run \ev{config/default.yaml:56-59}.

\paragraph{Why add a cross-encoder?} The baseline found the gold passage in the pool 96.7\% of the time but kept it in the top five only 88.0\% of the time, so reordering was the gap to close \ev{src/vidyarag/retrieve/rerank.py:8-12}. It raised recall@context to 0.913 and MRR to 0.830. \emph{Why MiniLM rather than \code{bge-reranker-base}?} 80\,MB against 1.04\,GB, on free hosting; the larger model was never measured, and the design notes say so. \emph{Why not Jina's reranker?} A non-commercial licence.

\paragraph{Why try query decomposition, and why reject it?} Reranking hurt multi-hop questions, apparently by favouring near-duplicates of the strongest match, so decomposition gave each hop its own search. It was rejected because multi-hop context recall fell from 0.861 to 0.778: rank fusion rewarded passages several sub-questions agreed on, which were the generic ones.

\paragraph{Why reciprocal rank fusion with $k = 60$?} Scores from different sub-queries are not comparable, so fusion uses ranks only; 60 is the published default, left untuned so as not to fit the evaluation \ev{src/vidyarag/retrieve/decompose.py:17-39}.

\paragraph{Why no hybrid (keyword plus dense) search?} Pool recall was already 0.967, and adding sparse vectors meant re-indexing \ev{src/vidyarag/store/collection.py:29-38}. \tagI The gold set does not test exact-term queries, where keyword search usually helps.

\section{Generation and the self-check}

\paragraph{Why Gemini?} Its free tier needs no card; the original OpenAI plan required billing \ev{docs/DESIGN.md}. \emph{Why the lite models?} \code{gemini-3.5-flash}'s free tier allowed 20 requests a day, measured; one evaluation needs about 60 generation calls.

\paragraph{Why pin exact model ids?} An alias such as \code{-latest} could change the model under a benchmark and make comparisons meaningless \ev{config/default.yaml:12-16}.

\paragraph{Why temperature 0?} For repeatability. It is not deterministic, which the project measured and accounts for with a noise band.

\paragraph{Why numbered citation markers and validation?} Models reproduce small numbers reliably but mangle identifiers, and \enquote{a fabricated citation is worse than none, because it is more convincing} \ev{src/vidyarag/generate/citations.py:7-23}.

\paragraph{Why a versioned prompt?} So every result can be attributed to the wording that produced it; the version is part of the trace and the cache key.

\paragraph{Why grade claims instead of the whole answer?} A single score gives a retry nothing to aim at; a failed claim becomes the next search \ev{src/vidyarag/correct/grader.py:7-12}.

\paragraph{Why a different model for grading?} A model grading its own output rates it favourably.

\paragraph{Why a refusal flag?} A refusal is perfectly supported by passages that lack the answer, so it scores 1.0; without the flag the loop would never abstain on the questions it exists to decline.

\paragraph{Why accept when grading fails?} Abstaining because of an API error \enquote{would be a silent, invisible failure --- the system would look appropriately cautious while actually being broken} \ev{src/vidyarag/correct/policy.py:64-69}. \emph{The honest follow-up:} it also means an answer can be returned unverified, and the demo does not currently say so.

\paragraph{Why thresholds of 0.8 and 0.5, untuned?} Fitting two thresholds to twelve unanswerable questions would look tuned, generalise no better, and use up the gold set as a development set \ev{src/vidyarag/correct/policy.py:9-18}.

\paragraph{Why at most two attempts?} Each attempt costs a generation and a grading call; \enquote{an unbounded loop is an unbounded bill and an unbounded latency.}

\section{Safety}

\paragraph{Why regular-expression guards?} The design constraint was false positives on textbook prose, measured against every passage; patterns are cheap, deterministic and explain themselves \ev{src/vidyarag/guard/patterns.py:22-28}. \emph{Why not a model-based classifier?} Not recorded. \tagI An extra model call per question on a rate-limited free tier, for a threat model where the corpus is trusted and the system prompt is public.

\paragraph{Why block questions but quarantine passages?} Sanitising an attack question would be \enquote{doing an attacker's editing for them}; refusing on a bad passage would let one document deny service to every question that retrieves it.

\section{Evaluation}

\paragraph{Why RAGAS?} Standard metrics are stronger evidence than hand-written ones a reader cannot inspect \ev{docs/DESIGN.md}. \emph{Why through an OpenAI-compatible client?} RAGAS detects asynchronous clients by an OpenAI attribute, which Google's own compatibility endpoint satisfies; no OpenAI account is involved.

\paragraph{Why deterministic retrieval metrics as well?} They are free, exact and locate faults that model-judged metrics cannot.

\paragraph{Why a language model to detect abstention?} The baseline refuses in varied prose, and keyword matching would be guesswork. \emph{The lesson:} the judge needed its own tests --- its output limit made it fail silently.

\paragraph{Why a frozen baseline and one change per profile?} So each difference can be attributed to one change.

\paragraph{Why withhold metrics above 10\% failures?} A run that lost 39 of 58 questions reported faithfulness 0.949 because the hard questions were missing.

\paragraph{Why cache answers?} On a free tier, a run that dies at question 19 should resume tomorrow rather than pay for those 19 again.

\section{Engineering}

\paragraph{Why no LangChain or LlamaIndex?} The pipeline is a few hundred lines over three libraries; a framework \enquote{would have added indirection without removing any of it.} LlamaIndex was planned, found unused, and removed.

\paragraph{Why FastAPI?} The README gives the practical reason: responses are Pydantic models, so the OpenAPI schema is generated from the same types the handlers return and cannot drift.

\paragraph{Why Gradio?} It is the SDK Hugging Face Spaces supports for this kind of demo, and it made showing the trace simple. \tagI for any comparison with Streamlit, which is not discussed.

\paragraph{Why Hugging Face Spaces on ZeroGPU without using a GPU?} ZeroGPU is the only free Gradio hardware; the platform requires a GPU-decorated function, so one is declared and never called.

\paragraph{Why uv?} So CI, the container and a developer's shell run identical commands from one lockfile \ev{Makefile}. \tagI Speed and built-in Python installation also help.

\paragraph{Why YAML profiles with \code{extra="forbid"}?} A profile is a diffable, citable description of one pipeline variant, and a typo fails at load instead of silently invalidating a benchmark.

\paragraph{Why structlog?} JSON logs are queryable in deployment, console output is readable locally \ev{src/vidyarag/observe/logging.py:1-5}.

\paragraph{Why a multi-stage Docker image without PyTorch?} About 400\,MB instead of about 2.5\,GB, enforced by a 900\,MB CI budget.

\paragraph{Why support Python 3.10?} The Space's platform pins 3.10 and ignores requests for another version.

\paragraph{Why mypy strict, ruff, black and gitleaks?} Type errors and style issues fail before merge; gitleaks runs first because \enquote{a leaked key in git history survives \code{git rm}.}
